\chapter{Curve-Shrinking Estimates}

This chapter develops the estimates needed to deform compact families of loops
by curve shrinking in a time-dependent Ricci-flow metric.  In particular, all
total-curvature and local regularity bounds retain their necessary dependence
on both initial total curvature and initial length.

\section{Curve-Flow Geometry}

Let $(M,g(t))$, $t\in I$, be a Ricci flow on a compact manifold, and let
$c\colon S^1\times I\to M$ be a smooth family of immersed curves satisfying
$\partial_t c=H=\nabla^g_SS$.  Write $X=c_*(\partial_x)$,
$S=X/|X|_{g(t)}$, and $k=|H|_{g(t)}$.

\begin{theorem}[Speed and tangent evolution]
\label{lem:curve-flow-speed-and-tangent-evolution}
\dcref{MT:ch19:19.6}
\source{morgan-tian-correction}
\notready
Let $\widehat c(x,t)=(c(x,t),t)$ and
$\widehat H=\widehat c_*(\partial_t)=H+\partial_t$.  Then
\[
 \widehat H(|X|^2)=-2\operatorname{Ric}_g(X,X)-2k^2|X|^2,
 \qquad
 [\widehat H,S]=(k^2+\operatorname{Ric}_g(S,S))S .
\]
Equivalently, the first formula is the time derivative of the squared speed in
the evolving pulled-back metric.
\end{theorem}
\begin{proof}
Because $X$ and $\widehat H$ are coordinate fields on the swept-out surface,
they commute.  Differentiating the metric and the vector field separately gives
\[
 \widehat H(|X|^2)=-2\operatorname{Ric}_g(X,X)
        +2\langle\nabla_X^gH,X\rangle .
\]
The fields $H$ and $X$ are orthogonal, while
$\nabla_XX=|X|^2\nabla_SS+X(|X|)S$; hence
$2\langle\nabla_XH,X\rangle=-2\langle H,\nabla_XX\rangle
=-2k^2|X|^2$.  This proves the first identity.  Finally,
\[
 [\widehat H,S]=[\widehat H,|X|^{-1}X]
 =-\frac{\widehat H(|X|^2)}{2|X|^2}S,
\]
and substitution of the first identity yields the second.
\end{proof}

\begin{lemma}[Space-time metric connection identities]
\label{lem:space-time-metric-connection}
\dcref{MT:ch19:19.7}
\source{morgan-tian-correction}
\notready
Equip $M\times I$ with the metric $\widehat g=g(t)+dt^2$, and denote its
Levi--Civita connection by $\nabla$.  If $A,B$ are horizontal vector fields,
then
\[
 \nabla_A B=\nabla_A^gB+\operatorname{Ric}_g(A,B)\partial_t,
 \qquad
 \langle\nabla_A\partial_t,B\rangle
       =-\operatorname{Ric}_g(A,B),
\]
and $\nabla_A\partial_t$ is horizontal.
\end{lemma}
\begin{proof}
In coordinates $x^1,\ldots,x^n,t$, the metric has components
$\widehat g_{ij}=g_{ij}(t)$, $\widehat g_{i0}=0$, and
$\widehat g_{00}=1$.  The Christoffel formula and
$\partial_tg_{ij}=-2\operatorname{Ric}_{ij}$ give
\[
 \widehat\Gamma^k_{ij}=\Gamma^k_{ij},\qquad
 \widehat\Gamma^k_{i0}=-g^{k\ell}\operatorname{Ric}_{i\ell},
 \qquad
 \widehat\Gamma^0_{ij}=\operatorname{Ric}_{ij},
\]
with all remaining symbols involving two time indices equal to zero.  Reading
off horizontal and vertical components gives all three assertions.
\end{proof}

\section{Product-Circle Geometry and Ramp Setup}

\begin{theorem}[Product-circle geometry]
\label{lem:product-circle-parallel-geometry}
\dcref{MT:ch19:19.11}
\source{morgan-tian}
\notready
On $(M,g(t))\times(S^1_\lambda,ds^2)$ let $U$ be the oriented unit tangent
to the circle factor.  Then $\nabla U=0$ and
$\operatorname{Ric}(V,U)=0$ for every tangent vector $V$; both assertions are
independent of $\lambda>0$.
\end{theorem}
\begin{proof}
The Levi--Civita connection of a Riemannian product is the product of the two
factor connections.  The unit tangent field on the flat one-dimensional
circle is parallel, hence so is its product lift.  The curvature tensor has no
mixed components and the circle Ricci tensor is zero, so contraction gives
$\operatorname{Ric}(V,U)=0$.
\end{proof}

\begin{definition}[Lifted ramp chart and circle branch]
\label{def:lifted-ramp-chart-branch}
\dcref{MT:ch19:19.12,MT:ch19:19.64}
\source{morgan-tian}
\notready
Throughout this chapter the parameter circle is
$S^1=\mathbb R/(2\pi\mathbb Z)$.  For every $\lambda>0$ regard
$S^1_\lambda=\mathbb R/(\lambda\mathbb Z)$ with the metric induced by the
Euclidean coordinate on the cover, so its circumference is $\lambda$, and
write $\pi_\lambda\colon M\times\mathbb R\to M\times S^1_\lambda$ for the
covering map induced by the universal cover of the circle factor.  The
degree-one map from the parameter circle to the target factor is therefore
$x\mapsto\lambda x/(2\pi)$ and has constant speed $\lambda/(2\pi)$.  A
degree-one ramp family is represented locally by a chosen
lift and a chosen lifted arc/branch over each polygonal side or graph interval.
All local charts use the fixed $M$ normal neighborhood of radius $r_0$ and an
affine coordinate on $\mathbb R$, then project by $\pi_\lambda$.  Hence chart
radii, connection bounds, and transition estimates are independent of
$\lambda$; changing the branch by a deck translation changes nothing.
\end{definition}
\begin{proof}
The product connection on $M\times\mathbb R$ is the lift of the product
connection on $M\times S^1_\lambda$, and deck translations are isometries.
Choose $r_0$ below the fixed normal radius in $M$ and use an interval in the
lifted $\mathbb R$ coordinate containing the selected arc.  Its projection is
injective on the chosen branch, and every other branch is an isometric deck
translate.  All local estimates therefore descend uniformly in $\lambda$.
\end{proof}

\begin{definition}[Ramp]
\label{def:ramp-curve}
\dcref{MT:ch19:19.12}
\source{morgan-tian}
\notready
For an immersed oriented curve $c\colon S^1\to M\times S^1_\lambda$ with
unit tangent $S$, put $u=\langle S,U\rangle$.  The curve is a \emph{ramp} if
$u(x)>0$ for every $x\in S^1$.
\end{definition}

\begin{theorem}[Evolution equation for the ramp slope]
\label{lem:ramp-slope-evolution}
\dcref{MT:ch19:19.11}
\source{morgan-tian}
\notready
Along a curve-shrinking flow in
$(M,g(t))\times(S^1_\lambda,ds^2)$, the function
$u=\langle S,U\rangle$ satisfies
\[
 \partial_tu=u''+(k^2+\operatorname{Ric}(S,S))u
 \ge u''-C'u,
\]
where $C'$ is any upper bound for the norm of the Ricci tensor of $g(t)$.
\end{theorem}
\begin{proof}
\uses{lem:curve-flow-speed-and-tangent-evolution,lem:product-circle-parallel-geometry}
The metric derivative contributes $-2\operatorname{Ric}(S,U)=0$, and $U$ is
parallel.  Thus the tangent commutator identity gives
\begin{align*}
 \partial_tu
 &=\langle\nabla_{\widehat H}S,U\rangle
 =\langle[\widehat H,S]+\nabla_S\widehat H,U\rangle\\
 &=(k^2+\operatorname{Ric}(S,S))u
   +S\langle H,U\rangle
 =(k^2+\operatorname{Ric}(S,S))u+S(Su).
\end{align*}
The lower bound follows from $k^2\ge0$ and
$\operatorname{Ric}(S,S)\ge-C'$.
\end{proof}

\begin{theorem}[Ramp preservation]
\label{thm:ramp-preservation}
\dcref{MT:ch19:19.13}
\source{morgan-tian}
\uses{def:ramp-curve}
\notready
Suppose $c$ is a curve-shrinking flow on $[t_0,T)$ in the ramp product and
$c(\cdot,t_0)$ is a ramp.  If $C'$ bounds $|\operatorname{Ric}|$, then
\[
 \min_{S^1}u(\cdot,t)\ge
 e^{-C'(t-t_0)}\min_{S^1}u(\cdot,t_0)>0
 \qquad(t_0\le t<T).
\]
In particular every $c(\cdot,t)$ is a ramp.
\end{theorem}
\begin{proof}
\uses{thm:time-slice-maximum-comparison,lem:ramp-slope-evolution}
Set $v=e^{C'(t-t_0)}u$.  The slope inequality in
\ref{lem:ramp-slope-evolution} says $\partial_tv\ge v''$.
Apply \ref{thm:time-slice-maximum-comparison} to $F=-v$ (with the
corresponding scalar upper differential inequality) on the closed circle.
This shows that $\min v(\cdot,t)$ is nondecreasing.  Evaluating it at $t_0$
and undoing the exponential factor gives the displayed estimate.
\end{proof}

\section{Graph-Flow Coordinate Estimates}

Fix a smooth curve-shrinking flow $c$ in a Ricci flow $(W,h(t))$ on
$[t_0,t_1]$, with initial total curvature and length at most $\Theta$ and
$L_0$.  Choose $0<\delta<1$ and $0<r_0\le1$ sufficiently small in terms of
$\Theta,L_0$, and the fixed ambient flow.  Then choose $0<r\le r_0$ and
$t'\le t_1-\delta r^2$.  Assume that $L(c_{t'})\ge r$, that every
length-$r$ subarc of $c_{t'}$ has total curvature at most $\delta$, and let
$t_2$ be maximal in $[t',t'+\delta r^2]$ while
$k^2\le2/(t-t')$.  These are the local data used throughout this section.

\begin{definition}[Local graph-flow coordinates]
\label{def:curve-flow-local-graph-coordinates}
\dcref{MT:ch19:19.64}
\source{morgan-tian}
\uses{def:lifted-ramp-chart-branch}
\notready
At the midpoint $p$ of a length-$r$ arc, lift the evolving arc through the
$h(t')$-exponential local diffeomorphism from a ball in $T_pW$.  Give
$T_pW$ the Euclidean metric $h'=h(t')_p$ and take the $z$-axis tangent to the
lifted arc.  The choice of $r_0$ ensures that throughout the tangent-space
ball
\[
 |h(t)-h'|_{h'}<\delta,\qquad
 |\Gamma(A,B)|_{h'}<\delta|A|_{h'}|B|_{h'}.
\]
On a fixed interval $I'$ of $h'$-length $l\in[0.8r,r]$, write the evolving
arc, up to its maximal graph time $t_3\le t_2$, as
$\widetilde\gamma(z,t)=(z,f(z,t))$, where $|f_z|_{h'}\le1/10$.  Define
\[
 Z=\partial_z\widetilde\gamma=(1,f_z),\qquad
 Y=\partial_t\widetilde\gamma=(0,f_t).
\]
Thus $Y$ is $h'$-orthogonal to the $z$-axis and differs from
$\nabla_SS$ by a scalar multiple of $Z$.
\end{definition}
\begin{proof}
The ambient curvature bounds and smoothness of the fixed Ricci flow give a
uniform radius on which the exponential map is a local diffeomorphism and the
pullbacks of all $h(t)$ are $C^1$-close to $h'$.  Lift the selected midpoint
and arc to the chosen branch in $M\times\mathbb R$, and choose $r_0$ inside
that fixed $M$ normal radius.  This construction depends on a
conjugate-radius bound, not on the injectivity radius.  It is therefore uniform
for $W=M\times S^1_\lambda$: the circle factor is flat, deck translations are
isometries, and the product curvature bounds are those of $M$ alone, hence
independent of $\lambda$.  The provisional curvature
bound gives
\[
 |\gamma_t(x)-\gamma_{t'}(x)|_{h'}
 \le\int_{t'}^t\sqrt{1+\delta}\,k(x,\tau)d\tau
 \le4\sqrt{t-t'}\le4\sqrt\delta r,
\]
so the central arc stays in the chart.  Its initial total turning is at most
$\delta$; parallel transport of its tangent, together with metric closeness,
makes the $z$ projection nonsingular and gives an initial graph with
$|f_z|_{h'}<2\delta$.  Continuity produces a maximal time $t_3\le t_2$ on
which the slope is at most $1/10$.  Removing small endpoint portions leaves a
fixed interval $I'$ of length between $0.8r$ and $r$.  Reparametrizing each
graph by its $z$ coordinate makes the time velocity vertical for $h'$, while
reparametrization changes geometric curve shrinking only by a tangential
multiple of $Z$.  This constructs precisely the declared fields $Z$ and $Y$.
\end{proof}

\begin{theorem}[Graph-flow velocity identities]
\label{lem:graph-flow-velocity-identities}
\dcref{MT:ch19:19.64}
\source{morgan-tian}
\uses{def:curve-flow-local-graph-coordinates}
\notready
On the graph interval,
\[
 Y=\frac{\nabla_ZZ-\langle\Gamma(Z,Z),\partial_z\rangle_{h'}Z}{|Z|^2}
 =\nabla_SS+\left(
 \frac{\langle\nabla_ZZ,Z\rangle}{|Z|^4}
 -\frac{\langle\Gamma(Z,Z),\partial_z\rangle_{h'}}{|Z|^2}
 \right)Z.
\]
\end{theorem}
\begin{proof}
In the fixed coordinates,
$\nabla_ZZ=(0,f_{zz})+\Gamma(Z,Z)$ and
$\langle Z,\partial_z\rangle_{h'}=1$.  Hence the first displayed
right-hand side is $h'$-orthogonal to the $z$-axis and differs from
$\nabla_ZZ/|Z|^2$ by a multiple of $Z$.  These two properties uniquely
characterize the graph velocity, proving the first equality.  The normal
projection formula
\[
 \nabla_SS=\frac{\nabla_ZZ}{|Z|^2}
 -\frac{\langle\nabla_ZZ,Z\rangle}{|Z|^4}Z
\]
then gives the second.
\end{proof}

\begin{definition}[Graph-flow connection drift]
\label{def:graph-flow-connection-drift}
\dcref{MT:ch19:19.65}
\source{morgan-tian}
\uses{def:curve-flow-local-graph-coordinates}
\notready
Define
\[
 \psi(Z)=\frac{\langle\Gamma(Z,Z),\partial_z\rangle_{h'}}{|Z|^2}.
\]
The metric and slope comparisons imply $|\psi(Z)|<1.5\delta$.
\end{definition}
\begin{proof}
The chart bound gives
$|\Gamma(Z,Z)|_{h'}<\delta|Z|_{h'}^2$.  Moreover
$|Z|_{h'}^2=1+|f_z|_{h'}^2$ and metric comparison gives
$|Z|^2\ge(1-\delta)|Z|_{h'}^2$.  Consequently
\[
 |\psi(Z)|
 \le\frac{\delta|Z|_{h'}^2}{(1-\delta)|Z|_{h'}^2}
 =\frac\delta{1-\delta}<1.5\delta
\]
after requiring $\delta<1/3$.
\end{proof}

\begin{theorem}[Small-slope graph and curvature comparisons]
\label{lem:graph-flow-small-slope-comparison}
\dcref{MT:ch19:19.65}
\source{morgan-tian}
\uses{def:curve-flow-local-graph-coordinates}
\notready
For sufficiently small $\delta$, wherever $Y\ne0$ the $h(t)$-angle between
$Y$ and $Z$ is greater than $\pi/4$, and
\begin{align*}
 &k\le|Y|\le\sqrt2k,\\
 &|\langle\nabla_ZZ,Z\rangle|<(k+2\delta)|Z|^3,\\
 &|\nabla_ZZ|<2(|Y|+\delta),\\
 &|\langle Y,Z\rangle|\le|Y|\,|f_z|_{h'}(1+3\delta).
\end{align*}
\end{theorem}
\begin{proof}
\uses{lem:graph-flow-velocity-identities,def:graph-flow-connection-drift}
Since $Y=(0,f_t)$ is $h'$-orthogonal to the $z$-axis and
$Z=(1,f_z)$ with $|f_z|_{h'}\le1/10$, their Euclidean angle is uniformly
separated from zero; $C^1$ metric closeness gives the stated $h(t)$ angle.
The normal component of $Y$ is
$\nabla_SS$, so the angle comparison yields
$k\le|Y|\le\sqrt2k$.  From
$\nabla_ZZ/|Z|^2=Y+\psi(Z)Z$, the bounds for $|Z|$ and
$|\psi|<1.5\delta$ give the third inequality.  The $h'$-orthogonality of
$Y$ to the $z$-axis gives
$|\langle Y,Z\rangle_{h'}|\le|Y|_{h'}|f_z|_{h'}$; metric comparison gives
the fourth.  Finally, taking the inner product of
$\nabla_ZZ/|Z|^2=Y+\psi Z$ with $Z$ and using the first and fourth bounds
gives the second inequality after decreasing $\delta$.
\end{proof}

\begin{theorem}[Derivatives of the graph-flow connection drift]
\label{lem:graph-flow-drift-derivative-bounds}
\dcref{MT:ch19:19.66}
\source{morgan-tian}
\uses{def:curve-flow-local-graph-coordinates}
\notready
In the fixed exponential chart write $\Gamma(t)$ for the Christoffel symbols
of $h(t)$ and assume $|\partial_t\Gamma(t)|_{h'}\le C_{\Gamma}$ on the
chart.  Thus the time derivative below includes the evolving-connection term.
For sufficiently small $\delta$ there is an ambient constant $C'$ such that
\[
 |Z(\psi(Z))|<C'(1+\delta|Y|),
 \qquad
 |Y(\psi(Z))|<C'(1+|Y|+\delta|\nabla_ZY|).
\]
\end{theorem}
\begin{proof}
\uses{lem:graph-flow-small-slope-comparison,def:graph-flow-connection-drift}
Differentiate the quotient defining $\psi$.  A $Z$ derivative falling on
$\Gamma$ is bounded by the ambient $C^1$ connection bound; one falling on a
$Z$ factor is bounded by $C\delta|\nabla_ZZ|$, and derivatives of the
denominator have the same form.  The estimate in
\ref{lem:graph-flow-small-slope-comparison} proves the first inequality.
Similarly, the evolving connection contributes
$\partial_t\Gamma(Z,Z)$, bounded by $C_{\Gamma}|Z|^2$, while the spatial
$Y$-derivative contributes at most $C'|Y|$ and derivatives of a
$Z$ factor contribute at most $C\delta|\nabla_YZ|$.  Since $Y$ and $Z$ are
coordinate fields, $\nabla_YZ=\nabla_ZY$, proving the second inequality.
\end{proof}

\section{Short-Time Existence}

\begin{lemma}[Strongly parabolic quasilinear equation on a closed circle]
\label{lem:strongly-parabolic-circle-equation}
\dcref{MT:ch19:19.1}
\source{topping:page-0058}
\source{topping:page-0059}
\notready
Let $E\to S^1$ be a smooth finite-dimensional vector bundle and let
$P\colon\Gamma(E)\to\Gamma(E)$ be a smooth second-order quasilinear operator.
Suppose that, on a neighbourhood of a section $w$, the linearisation has
principal symbol satisfying
\[
 \langle\sigma(DP(v))(x,\xi)z,z\rangle
 \ge \lambda_0|\xi|^2|z|^2
 \qquad(\lambda_0>0)
\]
for every $v$ in that neighbourhood, $x\in S^1$, and $\xi,z$ in the
corresponding fibres.  Then the initial-value problem
\[
 \partial_t v=P(v),\qquad v(0)=w,
\]
has a unique smooth solution on a positive time interval.  On a compact set
of initial sections on which the same symbol and coefficient bounds hold, the
interval can be chosen uniformly.  Solutions depend continuously on their
initial sections in $C^{2,\alpha}$ on every shorter interval, and a solution
whose ellipticity, speed, and $C^2$ bounds stay uniform at a finite endpoint
extends across that endpoint.
\end{lemma}
\begin{proof}
Topping's definition of the principal symbol and of a quasilinear operator on
pages~58--59 is exactly the displayed hypothesis; the same pages give
smooth short-time existence and uniqueness for a strongly parabolic equation
on a closed manifold.  For the uniform and continuous-dependence clauses, the
difference of two solutions satisfies a linear equation with leading
coefficient between positive constants.  Multiplying by successive spatial
derivatives and integrating over the periodic circle gives, for every $m$,
\[
 \frac{d}{dt}\|D^m v\|_{L^2}^2
 +\lambda_0\|D^{m+1}v\|_{L^2}^2
 \le C_m\sum_{j\le m}\|D^jv\|_{L^2}^2.
\]
The one-dimensional periodic Sobolev inequality and Gronwall's lemma convert
these estimates to $C^{2,\alpha}$ bounds and a difference estimate.  Compactness
of the initial set makes all constants and the first existence time common.
If the bounds persist up to a finite endpoint, the same estimates give a
convergent $C^{2,\alpha}$ endpoint section; applying the local existence
statement there extends the solution.  This is the parameter-uniform form of
the source theorem, with no target-shaped flow predicate.
\end{proof}

\begin{theorem}[Compact-parametric parabolic interface for immersed curve shrinking]
\label{thm:parametric-curve-shrinking-wellposedness}
\dcref{MT:ch19:19.1}
\source{morgan-tian}
\source{perelman-finite-extinction:page-0003}
\source{topping:page-0058}
\source{topping:page-0059}
\uses{lem:strongly-parabolic-circle-equation}
\notready
Morgan--Tian Claim~19.1 and Perelman's Section~2.1 statement give the
single-curve existence and curvature blow-up alternative.  The parameter-
uniform clauses below are obtained from the source-backed strongly parabolic
circle interface in \ref{lem:strongly-parabolic-circle-equation}; they are not
being attributed to the single-curve claim alone.
Let $P$ be compact, $0<\alpha<1$, and let $g(t)$ be a smooth Ricci flow on a
compact manifold for $t\in[t_0,t_1]$.  Suppose $p\mapsto c_{0,p}$ is
continuous in $C^{2,\alpha}$ and $|\partial_xc_{0,p}|\ge\iota_0>0$ for all
$p\in P$.  The gauged curve-shrinking equation is uniformly parabolic on a
time interval $[t_0,t_0+\tau]$, with $\tau>0$ depending only on the compact
initial set, $\iota_0$, and ambient $C^2$ bounds.  It has a unique solution in
the gauge and hence a unique geometric solution up to time-dependent
reparametrization; the solution depends continuously on $(p,t)$ in the
parabolic $C^{2,\alpha}$ topology.  If curvature and the immersion lower
bound remain uniformly bounded at a finite endpoint, the uniform periodic
parabolic estimates continue every member across that endpoint.  In particular a
compact parameter family has a common lifetime whenever these bounds are
uniform.
\end{theorem}
\begin{proof}
Fix the parameter $x$ on $S^1$ and replace geometric curve shrinking by the
gauged equation
\[
 \partial_tc=|c_x|_{g(t)}^{-2}\nabla_xc_x.                 \tag{G}
\]
Its normal component is $H=\nabla_SS$; the difference between its right-hand
side and $H$ is a scalar multiple of $c_x$.  Thus solutions of (G) and
geometric curve-shrinking flows differ only by time-dependent
reparametrization.  In a target chart, (G) has principal part
\[
 g^{xx}(c,c_x,t)\,\partial_x^2c,
 \qquad g^{xx}=|c_x|_{g(t)}^{-2},
\]
times the identity in the target coordinates.  On the set
$|c_x|\ge\iota_0/2$ this is a strictly parabolic quasilinear system.

The compact set $\{c_{0,p}:p\in P\}$ has one upper $C^{2,\alpha}$ bound, one
positive speed lower bound, and images contained in finitely many target
charts with uniformly bounded transition maps.  Equivalently, embed $M$
isometrically in Euclidean space, work in one fixed tubular neighborhood,
and project the lower-order terms back to $TM$.  Apply
\ref{lem:strongly-parabolic-circle-equation} to the section equation (G).
Its symbol is $|c_x|^{-2}|\xi|^2$ times the identity, so the common speed
lower bound gives the common ellipticity constant.  The difference estimate
in the proof of that lemma gives uniqueness in the gauge and continuous
dependence on $p$ in parabolic $C^{2,\alpha}$; differentiating (G) then
bootstraps smooth data.

For geometric uniqueness, solve the scalar ODE for the domain
diffeomorphisms whose velocity cancels the tangential part of (G).  The speed
lower bound makes this ODE smooth; after applying it, two geometric solutions
with the same initial curve solve (G), so gauged uniqueness identifies them.
At a finite endpoint, a common curvature bound and a common positive speed
lower bound first give common upper and lower length bounds by the speed
identity.  Reparametrize each curve at each time by constant speed, anchored
at the transported marked point $x=0$.  In this gauge the tangential second
derivative vanishes and
$\nabla_xc_x=|c_x|^2H$, so the curvature bound supplies a uniform $C^2$
bound for the parametrized curves.  The periodic energy estimates in
\ref{lem:strongly-parabolic-circle-equation} give uniform $C^{2,\alpha}$
bounds as $t\uparrow T$ (the compact initial strip is already controlled).
The solutions converge in $C^{2,\alpha}$ to a compact family of immersed
endpoint curves.  Applying the same uniform local construction at $T$
continues the family and completes the proof.
\end{proof}

\begin{theorem}[Short-time existence and curvature blow-up alternative]
\label{thm:curve-shrinking-short-time-existence}
\dcref{MT:ch19:19.1}
\source{morgan-tian}
\source{perelman-finite-extinction:page-0003}
\notready
Let $0<\alpha<1$ and let
$c_0\colon S^1\to(M,g(t_0))$ be an immersed $C^{2,\alpha}$ curve.  There is
$T>t_0$ and an immersed curve-shrinking flow $c$ on $[t_0,T)$.  If
$T<\sup I$
and $\sup_{S^1\times[t_0,T)}k<\infty$, then the flow extends, through
immersed curves, to an interval $[t_0,T+\eta)$ for some $\eta>0$.
Consequently every finite nonextendible endpoint has
$\limsup_{t\uparrow T}\max_{S^1}k(\cdot,t)=\infty$.
\end{theorem}
\begin{proof}
\uses{thm:parametric-curve-shrinking-wellposedness,
  lem:curve-flow-speed-and-tangent-evolution}
Apply \ref{thm:parametric-curve-shrinking-wellposedness} with a one-point
parameter space.  If $k$ stays bounded at a finite endpoint, the speed
identity in \ref{lem:curve-flow-speed-and-tangent-evolution} gives
$\partial_t\log|c_x|=-\operatorname{Ric}(S,S)-k^2$.  Direct integration on
$[t_0,t)$, using the assumed curvature bound and the bounded ambient Ricci
tensor, keeps $|c_x|$ bounded above and away from zero uniformly as
$t\uparrow T$.
The continuation clause of
\ref{thm:parametric-curve-shrinking-wellposedness} therefore extends the flow.
Taking the contrapositive proves the blow-up alternative.
\end{proof}

\section{Polygonal Ramp Approximation}

Let $\Lambda M$ denote the free loop space of $C^1$ maps $S^1\to M$, with
the $C^1$ topology.

For each $n\ge1$, let $\Lambda_n^{\mathrm{pg}}M$ denote the space of
$n$-sided piecewise-geodesic loops with the fixed equal subdivision of $S^1$,
topologized by their ordered vertices.  These polygonal spaces are auxiliary
topological targets; $\Lambda M$ itself still means the $C^1$ free loop space.

For a null-homotopic loop $c$, write $A(c)$ for the infimum of the areas of
piecewise smooth disks spanning $c$.

\begin{definition}[Regular polygonal approximation]
\label{def:regular-polygonal-loop-approximation}
\dcref{MT:ch19:19.18}
\source{morgan-tian}
\notready
Let $c\colon S^1\to M$ be $C^1$, let $n\ge1$, and put
$\xi_n=e^{2\pi i/n}$ and $p_j=c(\xi_n^j)$.  Choose a minimizing geodesic
$A_j$ from $p_j$ to $p_{j+1}$, with circular indexing, and parametrize it at
constant speed on the corresponding subinterval of $S^1$.  The resulting
piecewise geodesic loop $c_n$ is a \emph{regular $n$-polygonal approximation}
of $c$.  When all sides lie below the injectivity radius, the choices are
unique.
\end{definition}

\begin{theorem}[Polygonal approximation]
\label{lem:single-loop-polygonal-approximation}
\dcref{MT:ch19:19.19}
\source{morgan-tian}
\uses{def:regular-polygonal-loop-approximation}
\notready
For every $C^1$ loop $c\colon S^1\to M$ and every $\zeta>0$, there is
$N$ such that, for all $n\ge N$,
\[
 |L(c_n)-L(c)|<\zeta,
\]
and $c_n$ and $c$ cobound a piecewise smooth annulus of area less than
$\zeta$.
\end{theorem}
\begin{proof}
Uniform continuity of $c'$ in finitely many coordinate charts implies that
the sum of distances between successive sampled points converges to
$\int_{S^1}|c'|=L(c)$; this proves the length assertion.  For large $n$ the
curves $c$ and $c_n$ are uniformly closer than the injectivity radius.  Join
$c(x)$ to $c_n(x)$ by the unique short geodesic.  On each polygonal interval
this gives a smooth ruled strip, and the strips agree along their boundary
geodesics, producing a piecewise smooth annulus.  Its area is bounded by its
maximum transverse length times a uniform bound for the sum of longitudinal
lengths.  The transverse length tends uniformly to zero while the
longitudinal lengths converge to $L(c)$, so the area is less than $\zeta$
for all sufficiently large $n$.
\end{proof}

\begin{theorem}[Uniform polygonal approximation on compact families]
\label{lem:compact-family-polygonal-approximation}
\dcref{MT:ch19:19.20}
\source{morgan-tian}
\notready
Let $K$ be a compact set of $C^1$ loops, with the $C^1$ topology, and let
$\zeta>0$.  There is $N=N(K,\zeta)$ such that the two conclusions of
\ref{lem:single-loop-polygonal-approximation} hold for every $c\in K$ and
every $n\ge N$.
\end{theorem}
\begin{proof}
\uses{lem:single-loop-polygonal-approximation}
The proof of \ref{lem:single-loop-polygonal-approximation} depends only on the
modulus of continuity
of $c'$, a bound for $|c'|$, and a positive lower bound for the injectivity
radius on the common compact image.  Compactness of $K$ in the $C^1$ topology
gives one bound for $|c'|$ and one modulus of continuity for all $c\in K$:
otherwise sequences $c_j\in K$ and pairs of parameters converging together
would, after passage to a $C^1$-convergent subsequence, contradict continuity
of the derivative of the limit.  The union of the images of the loops is
compact, so its injectivity radius has a positive lower bound.  Applying the
estimates in \ref{lem:single-loop-polygonal-approximation} with these three
uniform data gives one
integer $N$ for every $c\in K$ and every $n\ge N$.
\end{proof}

\begin{theorem}[Polygonal approximation of sphere families]
\label{cor:sphere-family-polygonal-approximation}
\dcref{MT:ch19:19.21}
\source{morgan-tian}
\notready
Let $\Gamma\colon S^2\to\Lambda M$ be a continuous $C^1$-loop family whose
members are null-homotopic.  For every $\zeta>0$ there is $N$ such that, for
$n\ge N$, the regular polygonal family
$\Gamma_n\colon S^2\to\Lambda_n^{\mathrm{pg}}M$ is continuous, every
$\Gamma_n(q)$ is null-homotopic, and, uniformly in $q\in S^2$,
\[
 |L(\Gamma_n(q))-L(\Gamma(q))|<\zeta,
 \qquad
 |A(\Gamma_n(q))-A(\Gamma(q))|<\zeta.
\]
\end{theorem}
\begin{proof}
\uses{lem:compact-family-polygonal-approximation}
The image of $\Gamma$ is compact, so choose $n$ by uniform polygonal
approximation and also so large that every side is shorter than the ambient
injectivity radius.  Uniqueness and smooth dependence of short minimizing
geodesics on their endpoints make $q\mapsto\Gamma_n(q)$ continuous.  The
small annulus between $\Gamma(q)$ and $\Gamma_n(q)$ proves that the latter is
null-homotopic.  If this annulus has area $a<\zeta$, attaching it to any disk
spanning either boundary gives a disk spanning the other, whence
$|A(\Gamma_n(q))-A(\Gamma(q))|\le a<\zeta$.  The length conclusion is the
uniform estimate already obtained.
\end{proof}

\begin{theorem}[Length and total curvature of smoothed polygonal ramps]
\label{lem:smoothed-polygonal-ramp-bounds}
\dcref{MT:ch19:19.22}
\source{morgan-tian}
\notready
For each $q$, choose the fixed lifted polygonal branches in
$M\times\mathbb R$ and obtain $\widetilde\Gamma(q)$ from the $n$-polygonal
loop $\Gamma_n(q)$ by the standard smoothing with a fixed, nonnegative,
symmetric reparameterization bump.  Throughout this chapter we use the fixed
convention $S^1=\mathbb R/(2\pi\mathbb Z)$.  For $0<\lambda<1$, project the
degree-one lift $x\mapsto(\widetilde\Gamma(q)(x),\lambda x/(2\pi))$ by
$\pi_\lambda\colon M\times\mathbb R\to M\times S^1_\lambda$ and call the
resulting ramp $\widetilde\Gamma(q)^\lambda$.  Then
\[
 L(\widetilde\Gamma(q)^\lambda)
 \le L(\Gamma(q))+\lambda,
 \qquad
 \int_{\widetilde\Gamma(q)^\lambda}k\,ds\le n\pi .
\]
\end{theorem}
\begin{proof}
Each lifted polygonal side is minimizing in the fixed $M$ normal chart, so its
length is no greater than that of
the corresponding subarc of $\Gamma(q)$.  Smoothing only reparametrizes the
polygonal loop; hence $L(\widetilde\Gamma(q))\le L(\Gamma(q))$.  If
$a(x)dx$ is the arclength element of the base loop, the ramp arclength is
$\sqrt{a(x)^2+(\lambda/2\pi)^2}\,dx$, which is at most
$(a(x)+\lambda/2\pi)dx$.  Integration proves the first inequality.  On a
polygonal side the ramp lies in the totally geodesic flat product of that
geodesic with $S^1_\lambda$.  In parallel coordinates $(u,v)$ its tangent is
$b(x)u+(\lambda/2\pi)v$, where $b$ increases from zero on the first half of
the side and decreases symmetrically on the second.  On each half the tangent
turns monotonically within one quadrant, through an angle less than
$\pi/2$.  Thus the total turning per side is at most $\pi$; summing over the
$n$ sides proves the curvature estimate.
\end{proof}

\begin{lemma}[Free-loop approximation preserves the $\pi_1$-orbit]
\label{lem:free-loop-approximation-preserves-pi1-orbit}
\dcref{MT:ch18:18.16,MT:ch19:19.17}
\source{morgan-tian}
\notready
Assume $M$ is compact, connected, and $\pi_2(M)=0$.  If a family
$\Gamma\colon S^2\to\Lambda_0M$ represents the free-loop orbit
$[\xi]$ under
$\pi_2(\Lambda_0M,*)\cong\pi_3(M,*)$, then every family
$\widetilde\Gamma$ produced by the source's uniform controlled-ramp
approximation represents the same orbit
$[\xi]$.  No choice of a based representative is required.
\end{lemma}

\begin{proof}
The source's goodapprox construction supplies a homotopy
$S^2\times[0,1]\to\Lambda_0M$ between $\Gamma$ and
$\widetilde\Gamma$.  By the free-loop-space identification, free homotopy
classes of such families are exactly $\pi_1(M)$-orbits in $\pi_3(M,*)$.
Changing the marked point or lift during the homotopy changes a based
representative only by that action, so the orbit is unchanged.  This is the
translation from the based goodapprox statement in Morgan--Tian,
energy1.tex, lines~2041--2058, to the free-family convention used here.
\end{proof}

\begin{theorem}[Uniform smooth approximation by controlled ramps]
\label{thm:uniform-controlled-ramp-approximation}
\dcref{MT:ch19:19.17}
\source{morgan-tian}
\uses{cor:sphere-family-polygonal-approximation,
  lem:smoothed-polygonal-ramp-bounds,def:lifted-ramp-chart-branch}
\notready
Let $\Gamma\colon S^2\to\Lambda M$ be a continuous family of
null-homotopic loops.  Let $0<\zeta<1$.  There is a continuous family
$\widetilde\Gamma\colon S^2\to\Lambda M$ such that
$[\widetilde\Gamma]=[\Gamma]$ as a free-loop homotopy class, every member is
smooth, and, for every
$0<\alpha<1$, the same family is continuous in the $C^{2,\alpha}$ topology.
The homotopy is a free-loop homotopy; no based class or fixed-basepoint
preservation is asserted.  In particular,
\[
 |L(\widetilde\Gamma(q))-L(\Gamma(q))|<\zeta,\qquad
 |A(\widetilde\Gamma(q))-A(\Gamma(q))|<\zeta.
\]
Moreover there is $C_0<\infty$, independent of $q$ and
$0<\lambda<1$, for which both the length and total curvature of every ramp
$\widetilde\Gamma(q)^\lambda$ are at most $C_0$.
The homotopy assertion is a family assertion: the annuli form a continuous
map $S^2\times[0,1]\to\Lambda M$.
\end{theorem}
\begin{proof}
\uses{lem:smoothed-polygonal-ramp-bounds,cor:sphere-family-polygonal-approximation}
Choose one $n$ for which
\ref{cor:sphere-family-polygonal-approximation} holds.  On each
side reparametrize the minimizing geodesic by the fixed smooth bump that is
positive inside the side and vanishes to infinite order at its endpoints.
Because the bump vanishes to infinite order at every vertex, the resulting
family $\widetilde\Gamma$ consists of smooth loops.  The minimizing geodesics
depend smoothly on their endpoints inside the fixed normal neighborhood, and
the reparametrization is fixed; compactness of $S^2$ therefore makes the
family continuous in every $C^{2,\alpha}$ topology.  This reparametrization
preserves the image, length, and least spanning area.
The small annuli vary continuously and give the stated free-loop homotopy;
no based or fixed-basepoint preservation is part of this construction.
The estimate in
\ref{lem:smoothed-polygonal-ramp-bounds} gives total curvature at most
$n\pi$ and length at most
$\sup_qL(\Gamma(q))+\zeta+1$.  Taking the maximum of these two finite
numbers gives $C_0$.
\end{proof}

\section{Evolution Identities and Global Bounds}

Primes in this section denote covariant differentiation with respect to
arclength.  The symbol $\perp$ denotes projection orthogonal to $S$.

\begin{theorem}[Evolution equation for squared curvature]
\label{thm:curve-flow-curvature-evolution}
\dcref{MT:ch19:19.7}
\source{morgan-tian-correction}
\notready
Let $\widehat{\operatorname{Rm}}$ be the curvature
tensor of $\widehat g$.  At every smooth point of the flow,
\begin{align*}
 \partial_t k^2={}&(k^2)''-2| (\nabla_SH)^\perp|^2+2k^4
 -2\operatorname{Ric}_g(H,H)+4k^2\operatorname{Ric}_g(S,S)\\
 &+2\widehat{\operatorname{Rm}}(\widehat H,S,H,S)
 +2\operatorname{Ric}_g(S,S)\operatorname{Ric}_g(H,H)
 -2(\nabla^g_S\operatorname{Ric}_g)(S,H).
\end{align*}
\end{theorem}
\begin{proof}
\uses{lem:space-time-metric-connection,
  lem:curve-flow-speed-and-tangent-evolution}
By \ref{lem:space-time-metric-connection},
$\nabla_SS=H+\operatorname{Ric}_g(S,S)\partial_t$.  The speed and commutator
identities in \ref{lem:curve-flow-speed-and-tangent-evolution} give
\[
 \widehat H(|X|^2)=-2\operatorname{Ric}_g(X,X)-2k^2|X|^2.
\]
and $[\widehat H,S]=(k^2+\operatorname{Ric}_g(S,S))S$.  Now differentiate
$|H|^2$ in the $\widehat H$ direction and commute the two covariant
derivatives to obtain
\begin{align*}
 \partial_tk^2={}&2\langle\nabla_S\nabla_S\widehat H,H\rangle
 +2\langle\nabla_S[\widehat H,S],H\rangle
 +2\langle\nabla_{[\widehat H,S]}S,H\rangle\\
 &+2\widehat{\operatorname{Rm}}(\widehat H,S,H,S)
 -2\langle\nabla_{\widehat H}
  (\operatorname{Ric}_g(S,S)\partial_t),H\rangle .
\end{align*}
The last term equals
$2\operatorname{Ric}_g(S,S)\operatorname{Ric}_g(H,H)$.  Splitting
$\widehat H=H+\partial_t$ in the first term, the mixed connection identity
gives
\[
 2\langle\nabla_S\nabla_S\partial_t,H\rangle
 =-2(\nabla^g_S\operatorname{Ric}_g)(S,H)
  -2\operatorname{Ric}_g(H,H).
\]
Moreover,
\[
 2\langle\nabla_S\nabla_SH,H\rangle
 =(k^2)''-2|(\nabla_SH)^\perp|^2-2k^4,
\]
because $\langle\nabla_SH,S\rangle=-k^2$.  Each of the two commutator
terms contributes
$2(k^2+\operatorname{Ric}_g(S,S))k^2$.  Adding these displayed
identities gives exactly the asserted formula.
\end{proof}

\begin{theorem}[Regularized curvature inequality]
\label{lem:regularized-curve-curvature-inequality}
\dcref{MT:ch19:19.8}
\source{morgan-tian-correction}
\notready
There is a constant $C_1<\infty$, depending only on bounds for the ambient
Ricci flow and its first covariant curvature derivative, such that, for every
$\epsilon>0$ and
$h_\epsilon=(k^2+\epsilon^2)^{1/2}$,
\[
 \partial_t h_\epsilon\le h_\epsilon''+k^3+C_1(h_\epsilon+1).
\]
\end{theorem}
\begin{proof}
\uses{thm:curve-flow-curvature-evolution}
The final five terms of \ref{thm:curve-flow-curvature-evolution} are bounded in
absolute value by $\widehat C(k^2+k)$ for an ambient constant
$\widehat C$.  Since $h_\epsilon^2=k^2+\epsilon^2$,
\[
 \partial_t(h_\epsilon^2)
 \le (h_\epsilon^2)''-2|(\nabla_SH)^\perp|^2
       +2k^4+\widehat C(k^2+k).
\]
Now
$h_\epsilon'=\langle(\nabla_SH)^\perp,H\rangle/h_\epsilon$, so
$(h_\epsilon')^2\le |(\nabla_SH)^\perp|^2$.  Expanding
$(h_\epsilon^2)''$ and cancelling this nonpositive term yields
\[
 2h_\epsilon\partial_t h_\epsilon
 \le2h_\epsilon h_\epsilon''+2k^4+\widehat C(k^2+k).
\]
Divide by $2h_\epsilon$ and use $0\le k<h_\epsilon$; the result follows
with $C_1=\widehat C/2$.
\end{proof}

\begin{theorem}[Integral evolution inequalities]
\label{lem:curve-flow-length-total-curvature-evolution}
\dcref{MT:ch19:19.9}
\source{morgan-tian-correction}
\notready
Let
$L(t)=\int_{c_t}ds$ and $\Theta(t)=\int_{c_t}k\,ds$, and let $C_2$ bound
$|\operatorname{Ric}_g|$ on $M\times I$.  The function $L$ is smooth and
$\Theta$ is locally absolutely continuous.  For almost every time,
\[
 L'(t)\le\int_{c_t}(C_2-k^2)\,ds\le C_2L(t),
 \qquad
 \Theta'(t)\le(C_1+C_2)\Theta(t)+C_1L(t).
\]
\end{theorem}
\begin{proof}
\uses{lem:regularized-curve-curvature-inequality,lem:curve-flow-speed-and-tangent-evolution}
The speed identity in \ref{lem:curve-flow-speed-and-tangent-evolution} gives
\[
 \partial_tds=-(\operatorname{Ric}_g(S,S)+k^2)ds,
\]
which proves the length estimate.  Put
$\Theta_\epsilon=\int h_\epsilon ds$.  Differentiating and applying the
regularized curvature inequality gives
\begin{align*}
 \Theta_\epsilon'
 &\le\int\bigl(h_\epsilon''+k^3+C_1(h_\epsilon+1)
 -h_\epsilon(\operatorname{Ric}_g(S,S)+k^2)\bigr)ds\\
 &\le(C_1+C_2)\Theta_\epsilon+C_1L.
\end{align*}
Here $\int h_\epsilon''ds=0$ because the curve is closed, and
$k^3-h_\epsilon k^2\le0$.  Integrate the regularized inequality between two
times before sending $\epsilon\downarrow0$.  Monotone convergence gives the
corresponding integrated inequality for $\Theta$; local absolute continuity
then gives the displayed estimate almost everywhere.
\end{proof}

\begin{lemma}[Scalar exponential comparison]
\label{lem:coupled-length-curvature-comparison}
\dcref{MT:ch19:19.10}
\source{morgan-tian-correction}
\notready
Suppose $L,\Theta\colon[0,T]\to[0,\infty)$ are absolutely continuous and
\[
 L'\le C_2L,\qquad
 \Theta'\le(C_1+C_2)\Theta+C_1L
\]
almost everywhere.  Then, for $0\le t\le T$,
\[
 L(t)\le L(0)e^{C_2t},\qquad
 \Theta(t)+L(t)\le(\Theta(0)+L(0))e^{(C_1+C_2)t}.
\]
\end{lemma}
\begin{proof}
\uses{lem:curve-flow-length-total-curvature-evolution}
The derivative of $e^{-C_2t}L(t)$ is nonpositive almost everywhere, which
proves the first inequality.  Adding the two assumed inequalities and using
$C_1L+C_2L=(C_1+C_2)L$ gives
$(\Theta+L)'\le(C_1+C_2)(\Theta+L)$.  The same integrating-factor
argument proves the second inequality.
\end{proof}

\begin{theorem}[Exponential length and total-curvature bounds]
\label{thm:curve-flow-length-total-curvature-bounds}
\dcref{MT:ch19:19.10}
\source{morgan-tian-correction}
\notready
For the immersed curve-shrinking flow above and $t_0\le t\le t_1$,
\[
 L(t)\le L(t_0)e^{C_2(t-t_0)},
\]
and
\[
 \Theta(t)+L(t)
 \le(\Theta(t_0)+L(t_0))e^{(C_1+C_2)(t-t_0)},
 \qquad
 \Theta(t)\le(\Theta(t_0)+L(t_0))e^{(C_1+C_2)(t-t_0)}.
\]
Thus the total-curvature bound necessarily depends on both initial
total curvature and initial length.
\end{theorem}
\begin{proof}
\uses{lem:coupled-length-curvature-comparison}
Translate time by $t_0$ and apply
\ref{lem:coupled-length-curvature-comparison} to the inequalities for $L$ and
$\Theta$.  The first inequality in the second display uses only $L(t)\ge0$.
\end{proof}

\begin{theorem}[Global existence for ramps]
\label{thm:ramp-curve-shrinking-global-existence}
\dcref{MT:ch19:19.14}
\source{morgan-tian-correction}
\notready
Let $(M,g(t))$, $t_0\le t\le t_1$, be a Ricci flow and let
$c_0\colon S^1\to M\times S^1_\lambda$ be a $C^{2,\alpha}$ ramp at time
$t_0$, for some $0<\alpha<1$.
Then its curve-shrinking flow exists on all of $[t_0,t_1]$, remains immersed,
and consists entirely of ramps.
\end{theorem}
\begin{proof}
\uses{thm:time-slice-maximum-comparison,
  thm:parametric-curve-shrinking-wellposedness,
  thm:curve-shrinking-short-time-existence,
  lem:curve-flow-speed-and-tangent-evolution,
  lem:regularized-curve-curvature-inequality,
  lem:ramp-slope-evolution,thm:ramp-preservation}
Let $[t_0,T)$ be the maximal interval of existence.
Theorem~\ref{thm:ramp-preservation} gives
$u\ge a>0$ on every finite subinterval.  For $h=h_\epsilon$ compute, using
\ref{lem:regularized-curve-curvature-inequality} and the exact equation for
$u$ in \ref{lem:ramp-slope-evolution},
\[
 \partial_t(h/u)\le (h/u)''+2(u'/u)(h/u)'
             +C'(h+1)/u.
\]
Since $0<u\le1$ and $u\ge a$, the last term is at most
$C''(h/u+1)$.  At a spatial maximum the derivative terms have the favorable
sign, so the forward-difference form of
\ref{thm:time-slice-maximum-comparison} bounds
$\max(h/u)$ exponentially on $[t_0,T)$.  Letting $\epsilon\downarrow0$
bounds $k$ there.  The speed identity in
\ref{lem:curve-flow-speed-and-tangent-evolution} then gives uniform upper and
positive lower bounds for $|c_x|$ on $[t_0,T)$, exactly as in
\ref{thm:curve-shrinking-short-time-existence}.  If $T<t_1$, the blow-up
alternative in
\ref{thm:curve-shrinking-short-time-existence} extends the flow past $T$, a
contradiction.  Hence the flow reaches $t_1$; the endpoint clause of
\ref{thm:parametric-curve-shrinking-wellposedness} supplies its terminal
slice, and \ref{thm:ramp-preservation} supplies the last assertion.
\end{proof}

\section{Localized Total-Curvature Control}

Continue with the local data fixed at the start of the graph-flow section.
All constants below depend only on $\Theta$, $L_0$, and the fixed ambient
Ricci flow; decrease the already chosen $\delta$ and $r_0$ whenever a
displayed smallness condition requires it.

\begin{theorem}[Uniform bounds on evolved subarcs]
\label{lem:evolved-subarc-global-bounds}
\dcref{MT:ch19:19.59}
\source{morgan-tian-correction}
\notready
There is $D_0<\infty$, depending only on $\Theta$, $L_0$, and the fixed
ambient Ricci flow, such that every subarc $\gamma_{t'}$ of length at most
$r$ satisfies, for $t'\le t\le t'+\delta r^2$,
\[
 \int_{\gamma_t}k\,ds<D_0,\qquad L(\gamma_t)\le D_0r.
\]
\end{theorem}
\begin{proof}
\uses{thm:curve-flow-length-total-curvature-bounds}
The estimate in \ref{thm:curve-flow-length-total-curvature-bounds}, applied
to the entire closed curve, bounds its total curvature on the fixed interval
by a constant depending on
$\Theta,L_0$, and the ambient flow; restriction gives the first inequality.
For a transported subarc, the length evolution satisfies
$L'\le C_2L$, with no endpoint term because its parameter endpoints are
fixed.  Hence $L(\gamma_t)\le e^{C_2(t-t')}r$.  Enlarge one constant to
obtain both assertions.
\end{proof}

\begin{lemma}[Endpoint separation and transported-length lower bound]
\label{lem:transported-endpoint-separation}
\dcref{MT:ch19:19.60,MT:ch19:19.63}
\source{morgan-tian-correction}
\notready
Under the local data fixed above, every transported initial arc $E_{t'}$ of
length $r$ satisfies
\[
 L(E_t)\ge(1-\eta)r\qquad(t'\le t\le t_2),
\]
where $\eta\to0$ as $\delta\to0$, uniformly in the chosen arc.  Consequently,
for the cutoff chosen below to vanish when $0.45\le|s|/r\le0.5$, its support
stays a positive intrinsic distance (at least $0.04r$ after decreasing
$\delta$) from both endpoints of $E_t$.
\end{lemma}
\begin{proof}
\uses{lem:evolved-subarc-global-bounds}
The arclength identity gives
\[
 \frac d{dt}L(E_t)=-\int_{E_t}(\operatorname{Ric}(S,S)+k^2)\,ds.
\]
By \ref{lem:evolved-subarc-global-bounds},
$L(E_\tau)\le D_0r$ and $\int_{E_\tau}k\,ds\le D_0$.  On the provisional
interval the pointwise barrier gives
\[
 \int_{E_\tau}k^2\,ds
 \le \bigl(\max_{E_\tau}k\bigr)\int_{E_\tau}k\,ds
 \le D_0\sqrt{\frac2{\tau-t'}}.
\]
This is the integrable estimate used by Morgan--Tian; the nonintegrable bound
$k^2\le2/(\tau-t')$ alone would not suffice.  If $C_R$ bounds
$|\operatorname{Ric}|$, integration from $t'$ to
$t\le t'+\delta r^2$ gives
\[
 L(E_t)\ge r-C_RD_0r\,\delta r^2
       -2\sqrt2D_0\sqrt\delta\,r
 \ge(1-\eta)r,
\]
where $\eta=C_RD_0\delta r^2+2\sqrt2D_0\sqrt\delta$.  Choose $\delta$
uniformly so that $\eta<1/10$.  Apply the same calculation to each outer
portion, whose initial length is $0.05r$.  Its curvature term has the same
absolute $O(\sqrt\delta\,r)$ bound, so decreasing $\delta$ further makes its
evolved length at least $0.04r$.  This is the endpoint separation
needed for zero extension and integration by parts.
\end{proof}

Choose a length-$r$ arc centered at $x_0$, let
$s(x,t)=\int_{x_0}^x|X(y,t)|dy$, and fix a smooth cutoff $\chi$ supported in
$[-1/2,1/2]$, identically zero when $0.45\le |u|\le0.5$, equal to one on
$[-3/8,3/8]$, with
$|\chi'|+|\chi''|\le D'$.  Put $\varphi(x,t)=\chi(s(x,t)/r)$.

\begin{theorem}[Time derivative of the localized cutoff]
\label{lem:localized-cutoff-time-derivative}
\dcref{MT:ch19:19.60}
\source{morgan-tian-correction}
\notready
There is $D_1$, depending only on $\Theta,L_0$, and the ambient flow, such
that, on $[t',t_2]$,
\[
 |\partial_t\varphi(x,t)|
 \le D_1\left(1+\frac1{r\sqrt{t-t'}}\right).
\]
\end{theorem}
\begin{proof}
\uses{lem:evolved-subarc-global-bounds,lem:curve-flow-speed-and-tangent-evolution}
The speed evolution formula gives
\[
 |\partial_ts(x,t)|
 \le C L(\gamma_t)+\int_{x_0}^x k^2ds.
\]
The speed and curvature bounds imply
\[
 \int_{x_0}^xk^2ds
 \le\sqrt{\frac2{t-t'}}\int_{\gamma_t}kds
 \le\frac{\sqrt2D_0}{\sqrt{t-t'}}.
\]
Since $\partial_t\varphi=r^{-1}\chi'(s/r)\partial_ts$, these estimates give
$|\partial_t\varphi|\le D_1(1+(r\sqrt{t-t'})^{-1})$ after absorbing
$D'D_0$ and fixed ambient constants into $D_1$.
\end{proof}

\begin{theorem}[Localized total-curvature differential inequality]
\label{lem:localized-total-curvature-evolution}
\dcref{MT:ch19:19.61}
\source{morgan-tian-correction}
\notready
There are constants $D_2,D_3<\infty$, depending only on
$\Theta,L_0$, and the ambient flow, such that the weighted curvature is
locally absolutely continuous and, for almost every $t' < t\le t_2$,
\[
 \frac d{dt}\int_{\gamma_t}\varphi k\,ds
 \le D_2\left(1+\frac1{r\sqrt{t-t'}}\right)+\frac{D_2}{r^2}
       +D_3\int_{\gamma_t}\varphi k\,ds
\]
\end{theorem}
\begin{proof}
\uses{lem:regularized-curve-curvature-inequality,lem:evolved-subarc-global-bounds,lem:transported-endpoint-separation,lem:localized-cutoff-time-derivative,lem:curve-flow-length-total-curvature-evolution}
First apply \ref{lem:transported-endpoint-separation}; this is what makes the
zero extension of $\varphi$ smooth at the transported endpoints and licenses
the integrations by parts below.
Differentiate the regularized integral
$\int\varphi h_\epsilon ds$.  The term in which the derivative hits
$\varphi$ is bounded by \ref{lem:localized-cutoff-time-derivative} times
$\int_{\gamma_t}h_\epsilon ds$; after letting $\epsilon\downarrow0$, this
is at most $D_1D_0(1+(r\sqrt{t-t'})^{-1})$.  Combining
\ref{lem:regularized-curve-curvature-inequality} with the length-element
evolution in \ref{lem:curve-flow-length-total-curvature-evolution} leaves a
diffusion term, a multiple of the weighted curvature, and the length term
$C'_1\int\varphi ds$.  Extend $\varphi$ by zero to the closed curve.  Its
support stays away from the subarc endpoints, so two integrations by parts
give
\[
 \left|\int\varphi h_\epsilon''ds\right|
 =\left|\int\varphi''h_\epsilon ds\right|
 \le\frac{D'}{r^2}\int_{\gamma_t}h_\epsilon ds.
\]
Lemma~\ref{lem:evolved-subarc-global-bounds} bounds the last integral by
$D_0+\epsilon D_0r$, uniformly as $\epsilon\downarrow0$.  Finally
$C'_1\int\varphi ds\le C'_1L(\gamma_t)\le C'_1D_0r$ and is absorbed in
$D_2$.  Integrate the regularized inequality before letting
$\epsilon\downarrow0$.  Monotone convergence and the preceding uniform bounds
give the displayed almost-everywhere upper derivative after enlarging
$D_2,D_3$.  This is the direction needed below; the evolution inequality does
not assert a corresponding lower bound.
\end{proof}

\begin{theorem}[Small localized total-curvature bound]
\label{lem:localized-total-curvature-smallness}
\dcref{MT:ch19:19.62}
\source{morgan-tian-correction}
\notready
There is $D_4$, depending only on $\Theta,L_0$, and the ambient flow, such
that
\[
 \int_{\gamma_t}\varphi k\,ds\le D_4\sqrt\delta
 \qquad(t'\le t\le t_2).
\]
\end{theorem}
\begin{proof}
\uses{lem:localized-total-curvature-evolution}
Put $y(t)=\int\varphi kds$.  The estimate in
\ref{lem:localized-total-curvature-evolution} and the integrating
factor $e^{-D_3(t-t')}$ give
\[
 y(t)\le e^{D_3(t-t')}
 \left(y(t')+D_2\int_0^{t-t'}
   \left(1+\frac1{r\sqrt s}+\frac1{r^2}\right)ds\right).
\]
Here $y(t')\le\delta$ and $t-t'\le\delta r^2$.  The integral is at most
$\delta r^2+2\sqrt\delta+\delta$.  Since $r,\delta<1$, all terms are bounded
by a fixed multiple of $\sqrt\delta$, proving the assertion.
\end{proof}

\begin{theorem}[Short-subarc total-curvature bound]
\label{lem:short-subarc-total-curvature-bound}
\dcref{MT:ch19:19.63}
\source{morgan-tian}
\notready
For $t\in[t',t_2]$, every transported initial subarc of length at most $r$
has total curvature at most $2D_4\sqrt\delta$.  Every subarc of $c_t$ of
length at most $r/2$ has the same bound.
\end{theorem}
\begin{proof}
\uses{thm:curve-flow-length-total-curvature-bounds,lem:localized-total-curvature-smallness}
Divide an initial arc of length at most $r$ into two arcs of length at most
$r/2$.  Enlarge each to be the central region, where $\varphi=1$, of a
length-$r$ arc.  Applying
\ref{lem:localized-total-curvature-smallness} to both gives the first
assertion.  On the provisional interval the speed formula, the defining
provisional curvature bound, and
\ref{thm:curve-flow-length-total-curvature-bounds} give the needed backward
length comparison as follows.  For a transported arc $E_t$,
direct differentiation of its arclength gives
\[
 \frac d{dt}L(E_t)
 =-\int_{E_t}(\operatorname{Ric}(S,S)+k^2)ds.
\]
Indeed, the metric derivative contributes $-\operatorname{Ric}(S,S)ds$,
and differentiating the tangent under $\partial_tc=\nabla_SS$ contributes
$-k^2ds$.  Thus
\[
 L(E_{t'})-L(E_t)
 \le C(t-t')L(E_{t'})+
 \int_{t'}^t\sqrt{\frac2{\tau-t'}}\,\Theta(\tau)d\tau.
\]
The estimate in \ref{thm:curve-flow-length-total-curvature-bounds} bounds
$\Theta(\tau)$ uniformly, so the last term is $O(\sqrt\delta r)$.  When
$L(E_{t'})=r$, choosing $\delta$
relative to the fixed global bounds gives $L(E_t)\ge0.9r$.  If the preimage
at $t'$ of a time-$t$ arc $J$ of length at most $r/2$ had length greater
than $r$, it would contain an initial length-$r$ subarc $E_{t'}$; its evolved
image $E_t\subset J$ would have length at least $0.9r$, a contradiction.
Thus that preimage has length at most $r$, and the first assertion applies.
\end{proof}

\section{Short-Time Curvature Barrier}

\begin{theorem}[Full-interval graph-slope persistence]
\label{lem:graph-slope-persistence}
\dcref{MT:ch19:19.58}
\source{morgan-tian}
\uses{def:curve-flow-local-graph-coordinates}
\notready
For sufficiently small $\delta$, the maximal graph time equals the
provisional curvature time, $t_3=t_2$, and
\[
 |f_z(z,t)|_{h'}<\frac1{10}
 \qquad((z,t)\in I'\times[t',t_2]).
\]
\end{theorem}
\begin{proof}
\uses{lem:graph-flow-velocity-identities,lem:graph-flow-small-slope-comparison,lem:graph-flow-drift-derivative-bounds,lem:short-subarc-total-curvature-bound,def:graph-flow-connection-drift,thm:curve-flow-length-total-curvature-bounds}
Differentiating $|Z|^2$ and using
\ref{lem:graph-flow-velocity-identities} gives
\[
 \partial_t|Z|^2
 =Z\left(\frac{Z(|Z|^2)}{|Z|^2}\right)-2|Y|^2|Z|^2+V,
 \qquad |V|\le C'(1+\delta|Y|).
\]
Integrating over $I'=(0,l)$, discarding the nonpositive $Y$ term, and using
$|Y|\le\sqrt2k$ gives
\[
 \frac d{dt}\int_{I'}|Z|^2dz
 \le \left.\frac{Z(|Z|^2)}{|Z|^2}\right|_0^l
      +C'\int_{I'}(1+\delta k)dz.
\]
The second estimate in \ref{lem:graph-flow-small-slope-comparison} gives
\[
 \frac{Z(|Z|^2)}{|Z|^2}
 =\frac{2\langle\nabla_ZZ,Z\rangle}{|Z|^2},
\]
whose absolute value at either endpoint is at most $C(k+\delta)$.  Hence the
provisional bound makes its time integral at most
$C\sqrt\delta r$.  Metric comparison changes $dz$ to $ds$ by a fixed factor,
so \ref{thm:curve-flow-length-total-curvature-bounds} controls
$\int_{I'}k\,dz$ uniformly.  Since $t-t'\le\delta r^2$ and $r\le1$, the
remaining time integral is at most $C(\delta r^3+\delta^{3/2}r^2)$.
At $t'$ the initial slope bound in
\ref{def:curve-flow-local-graph-coordinates} gives
$\int_{I'}|Z|^2dz\le(1+3\delta)l$.  Subtracting the lower metric bound
$|Z|^2\ge(1-\delta)(1+|f_z|_{h'}^2)$ therefore yields
\[
 \int_{I'}|f_z|_{h'}^2dz
 \le4\delta l+C''(\sqrt\delta r+\delta r^3+\delta^{3/2}r^2).
\]
Since $0.8r\le l\le r$, some $z(t)\in I'$ has
$|f_z(z(t),t)|_{h'}\le1/20$ when $\delta$ is small.  The total turning of
$I'$ is at most $2D_4\sqrt\delta$ by
\ref{lem:short-subarc-total-curvature-bound}.  Parallel
transporting unit tangents along the arc shows that their angles differ by
at most this total turning plus the $C^1$ coordinate error.  Decreasing
$\delta$ once more therefore gives $|f_z|_{h'}<1/10$ everywhere.  This is a
strict improvement of the condition defining $t_3$; continuity and
maximality force $t_3=t_2$.
\end{proof}

\begin{theorem}[Tangential component of $\nabla_ZY$]
\label{lem:graph-flow-tangential-derivative}
\dcref{MT:ch19:19.67}
\source{morgan-tian}
\uses{def:curve-flow-local-graph-coordinates,lem:graph-slope-persistence}
\notready
On $I'\times[t',t_2]$,
\[
 |\langle\nabla_ZY,Z\rangle|
 \le\bigl(|f_z|_{h'}(|\nabla_ZY|+2\delta|Y|)
       +\delta|Z|^2|Y|\bigr)(1+\delta).
\]
\end{theorem}
\begin{proof}
Write $Y=(0,\phi)$.  Then
$\nabla_ZY=(0,\phi_z)+\Gamma(Z,Y)$, and hence
\[
 |\langle\nabla_ZY,Z\rangle_{h'}|
 \le |f_z|_{h'}|\phi_z|_{h'}+\delta|Z|^2|Y|.
\]
The same identity gives
$|\phi_z|_{h'}\le|\nabla_ZY|+\delta|Z||Y|$.  Using
$|Z|\le1+\delta$, then comparing $h'$ with $h(t)$, proves the displayed
bound.
\end{proof}

\begin{theorem}[Short-time pointwise curvature barrier]
\label{thm:short-time-curve-curvature-barrier}
\dcref{MT:ch19:19.58}
\source{morgan-tian}
\uses{def:graph-flow-connection-drift,
  thm:time-slice-maximum-comparison,lem:graph-flow-small-slope-comparison,
  lem:graph-flow-drift-derivative-bounds,lem:graph-flow-tangential-derivative,
  def:curve-flow-local-graph-coordinates,lem:graph-slope-persistence}
\notready
Given bounds $\Theta,L_0$ and the fixed ambient Ricci flow, there are
$\delta>0$ and $0<r_0\le1$ such that the localized hypotheses stated above
imply
\[
 k(x,t)^2\le\frac2{t-t'}
 \qquad(x\in S^1,\ t'<t\le t'+\delta r^2).
\]
The constants may depend on both $\Theta$ and $L_0$.
\end{theorem}
\begin{proof}
\uses{thm:time-slice-maximum-comparison,lem:graph-flow-small-slope-comparison,lem:graph-flow-drift-derivative-bounds,lem:graph-flow-tangential-derivative,def:curve-flow-local-graph-coordinates,lem:graph-slope-persistence}
Work on the maximal provisional interval $[t',t_2]$.
Lemma~\ref{lem:graph-slope-persistence} makes
\ref{lem:graph-flow-small-slope-comparison},
\ref{lem:graph-flow-drift-derivative-bounds}, and
\ref{lem:graph-flow-tangential-derivative} valid throughout it.  Put
$a=|Z|^2$, $b=|Y|^2$, $P=\nabla_ZY$, and $e=\langle Y,Z\rangle$.  Since
$Z$ and $Y$ are the coordinate fields of the graph map, they commute.  The
metric variation and torsion-freeness therefore give
\[
 a_t=-2\operatorname{Ric}(Z,Z)+2\langle\nabla_YZ,Z\rangle
     =-2\operatorname{Ric}(Z,Z)+2\langle\nabla_ZY,Z\rangle.
\]
Substitute $Y=\nabla_ZZ/a-\psi(Z)Z$.  Differentiating the quotient and
using $a_z=2\langle\nabla_ZZ,Z\rangle$ yields
\[
 a_t=Z\left(\frac{Za}{a}\right)-2ab+V_0,
\]
where
\[
 V_0=-4a\langle Y,\psi(Z)Z\rangle-2\psi(Z)^2a^2
 -2\langle\nabla_Z(\psi(Z)Z),Z\rangle-2\operatorname{Ric}(Z,Z).
\]
Now $a_z=2ae+2\psi(Z)a^2$, so expanding the first term and absorbing the
terms containing $\psi$ into $V$ gives
\begin{align*}
 a_t&=\frac{a_{zz}}a-2ab-4e^2+V.                    \tag{1}
\end{align*}
The bounds in \ref{def:graph-flow-connection-drift} and
\ref{lem:graph-flow-drift-derivative-bounds}, together with
\ref{lem:graph-flow-small-slope-comparison}, show that
\[
 |V|\le C'(1+\delta|Y|).                            \tag{2}
\]

For $b$, differentiate once more, commute $\nabla_Y$ and $\nabla_Z$, and
then use $Y=\nabla_ZZ/a-\psi(Z)Z$.  Before estimating any term this gives
\begin{align*}
 b_t={}&\frac{b_{zz}}a-\frac{2|P|^2}a
 -\frac4a\left\langle\frac{\nabla_ZZ}{a},Y\right\rangle
             \langle P,Z\rangle
 -2\operatorname{Ric}(Y,Y)\\
 &+\frac{2\operatorname{Rm}(Y,Z,Y,Z)}a
 -2\langle\nabla_Y(\psi(Z)Z),Y\rangle.
\end{align*}
Since
$\langle\nabla_ZZ/a,Y\rangle=b+\psi(Z)e$, the estimates in
\ref{def:graph-flow-connection-drift} and
\ref{lem:graph-flow-drift-derivative-bounds} reduce this identity to
\[
b_t=\frac{b_{zz}}a-\frac{2|P|^2}a
       -\frac{4b}{a}\langle P,Z\rangle+W,
 \qquad
 |W|\le C'|Y|\bigl(1+|Y|+\delta|P|\bigr).           \tag{3}
\]

Set $Q=b/(2-a)$.  The slope bound in
\ref{def:curve-flow-local-graph-coordinates} implies
$1-\delta\le a<1.01(1+\delta)$, so the denominator is positive.  Applying
the quotient rule to (1) and (3), and comparing the result with the second
$z$-derivative of $Q$, gives
\begin{align*}
 Q_t={}&\frac{Q_{zz}}a
 -\frac{8\langle P,Y\rangle e}{(2-a)^2}
 -\frac{8ba e^2}{(2-a)^3}
 -\frac{2|P|^2}{a(2-a)}
 -\frac{4b\langle P,Z\rangle}{a(2-a)}\\
 &-\frac{2ab^2}{(2-a)^2}
 -\frac{4be^2}{(2-a)^2}+A,                          \tag{4}
\end{align*}
where (2)--(3), $a_z=2ae+2\psi a^2$, and the bounds for $a$ imply
\[
|A|\le C'\bigl(\sqrt b+b+\delta |P|\sqrt b+\delta b^{3/2}\bigr).\tag{5}
\]
Indeed, the four terms in $Q_{zz}$ not containing $a_{zz}$ or $b_{zz}$
are obtained from
\[
 b_z=2\langle P,Y\rangle,
 \qquad a_z=2ae+2\psi a^2;
\]
expanding their product gives the first two terms after $Q_{zz}/a$ in
(4), while every remaining factor of $\psi$ satisfies (5).

The fourth estimate in \ref{lem:graph-flow-small-slope-comparison} gives
$|e|\le\sqrt b\,|f_z|_{h'}(1+3\delta)$, and
\ref{lem:graph-flow-tangential-derivative} controls
$|\langle P,Z\rangle|$.  Discard the two nonpositive $e^2$ terms in
(4), use $|f_z|_{h'}<1/10$, and apply
\[
\sqrt b\le 1+b,
 \qquad
\delta |P|\sqrt b+\delta b^{3/2}
 \le \delta\bigl(|P|^2+b^2+2b\bigr).
\]
After decreasing $\delta$, all terms of degree two in $|P|$ and $b$ are
bounded above by
\[
 1.6b|P|-0.8|P|^2-0.8b^2\le0.
\]
Thus, for an ambient constant $C_*>0$,
\[
 Q_t\le\frac{Q_{zz}}a-Q^2+C_*(1+Q).
\]
Choose $K\ge\max\{C_*,1\}$.  Since
$C_*(1+Q)\le2KQ+K^2$, completing the square gives
\[
 Q_t\le\frac{Q_{zz}}a-(Q-K)^2+2K^2.                            \tag{6}
\]

Let $I'=(0,l)$, where $0.8r\le l\le r$.  For
$g(z)=l^2/(z^2(l-z)^2)$ direct differentiation gives
$g_{zz}\le12g^2$.
Writing $\alpha=(1-\delta)^{-1}$, set the harmless barrier constant to $2K$:
\[
 B(z,t)=\frac1{t-t'}+12\alpha g(z)+2K.
\]
Then
\[
 -B_t+\alpha B_{zz}+2K^2
 \le\frac1{(t-t')^2}+144\alpha^2g^2+2K^2\le B^2.                 \tag{7}
\]
Put $F=Q-K-B$.  If $F$ has a nonnegative spatial maximum, then
$F_{zz}\le0$, $Q-K\ge B>0$, and (6)--(7), together with
$a\ge1-\delta$, imply
\[
 F_t\le\frac{F_{zz}}a+\alpha B_{zz}-B_t
       -(Q-K)^2+2K^2\le0.
\]
Moreover $F$ tends uniformly to $-\infty$ as $t\downarrow t'$ and tends
to $-\infty$ at both spatial endpoints.  Applying
\ref{thm:time-slice-maximum-comparison} therefore gives
$F\le0$ on $(0,l)\times(t',t_2]$.
At the midpoint,
\[
 Q(l/2,t)\le\frac1{t-t'}+
 \frac{192(1-\delta)^{-1}}{l^2}+3K
 <\frac{3}{2(t-t')},
\]
because $(t-t')/l^2\le\delta/(0.8)^2$ and
$t-t'\le\delta r_0^2$.  Here $\delta$ is chosen in terms of the ambient
constant $K$ and the earlier constants $D_i$, which depend on
$\Theta,L_0$, and the ambient flow.  Finally
$k^2\le b=(2-a)Q\le(1+\delta)Q<2/(t-t')$ after also imposing
$\delta<1/3$.  The initial midpoint was arbitrary, so the strict inequality
holds everywhere on $[t',t_2]$.  Compactness then extends the provisional
inequality past any endpoint smaller than $t'+\delta r^2$; maximality forces
$t_2=t'+\delta r^2$.
\end{proof}

\begin{lemma}[Parabolic regularity for graph curve flow]
\label{lem:curve-flow-parabolic-regularity-interface}
\dcref{MT:ch19:19.24}
\source{morgan-tian}
\source{morgan-tian-correction}
\source{topping:page-0050}
\source{topping:page-0058}
\source{topping:page-0059}
\uses{def:curve-flow-local-graph-coordinates,lem:graph-slope-persistence,
  lem:graph-flow-small-slope-comparison,lem:strongly-parabolic-circle-equation}
\notready
On a graph interval on which the slope bound gives
$0<a_-\le |Z|^2\le a_+<\infty$, suppose the ambient metric and connection
coefficients have bounded derivatives of every order and the curve-flow
curvature barrier gives $|H|^2\le C_0(t-t')^{-1}$.  Then, for every $m\ge1$,
the uniformly parabolic gauged equation has interior (and periodic-boundary)
estimates
\[
 |\nabla_S^mH|^2\le C_m(t-t')^{-(m+1)},
\]
where $C_m$ depends only on $a_\pm$, the ambient coefficient bounds, and
$C_0$.
\end{lemma}
\begin{proof}
Morgan--Tian Lemma~19.24 supplies the $m=0$ barrier and the corrected
dependence on initial length and total curvature.  In the graph coordinates,
the principal coefficient is $a^{-1}$, uniformly elliptic by the slope
hypothesis, and the symbol calculation is the one in
\ref{lem:strongly-parabolic-circle-equation}.  Commuting $m$ spatial
derivatives with the gauged equation gives a linear uniformly parabolic
equation for $\nabla_S^mH$; every remaining term is a polynomial in lower
derivatives of $H$ and in the bounded ambient coefficients.  On a periodic
subinterval, integration by parts gives
\[
 \frac d{dt}\|D^mH\|_2^2+a_-\|D^{m+1}H\|_2^2
 \le C_m\left(\|D^mH\|_2^2+
       \sum_{j<m}\|D^jH\|_\infty^2\right).
\]
The one-dimensional Sobolev inequality and parabolic rescaling by
$(z,t-t')\mapsto((t-t')^{-1/2}z,1)$ turn the induction hypothesis into the
next $L^\infty$ bound.  Starting with the curvature barrier at $m=0$ and
inducting gives
$|\nabla_S^mH|^2\le C_m(t-t')^{-(m+1)}$.  Topping's discussion of the same
higher-derivative induction (page~50) and its strongly parabolic symbol and
existence interface (pages~58--59) provide the registered analytic inputs;
no unregistered Altschuler theorem is required.
\end{proof}

\begin{theorem}[Derivative estimates for curve shrinking]
\label{thm:curve-shrinking-derivative-estimates}
\dcref{MT:ch19:19.24}
\source{morgan-tian}
\source{morgan-tian-correction}
\source{topping:page-0050}
\source{topping:page-0058}
\source{topping:page-0059}
\uses{thm:short-time-curve-curvature-barrier,
  def:curve-flow-local-graph-coordinates,lem:graph-slope-persistence,
  lem:graph-flow-small-slope-comparison,
  lem:curve-flow-parabolic-regularity-interface}
\notready
The correction supplies the dependence of the constants in Lemma~19.24 on
the initial length and total-curvature bounds.  The all-order estimate is the
explicit commuted-equation induction in
\ref{lem:curve-flow-parabolic-regularity-interface}, rather than an
unstated appeal to a separate theorem.
Let $P$ be compact and let $c_p$, $p\in P$, be a family of the gauged curve
flows covered by \ref{thm:short-time-curve-curvature-barrier}, with common
initial length bound $L_0$, total-curvature bound $\Theta_0$, and the common
small-turning constants of that theorem.  The ambient flow may be
$M\times S^1_\lambda$ for any $0<\lambda<1$; its curvature and connection
bounds are those of $M$ and are independent of $\lambda$.  There are constants
$\widetilde C_i<\infty$, depending on the ambient Ricci flow and on the
uniform initial total-curvature and length bounds, but independent of
$p\in P$ and $0<\lambda<1$, such that
for $t'<t<t'+\delta r^2$,
\[
 k^2\le \widetilde C_0(t-t')^{-1},\qquad
 |\nabla_SH|^2\le\widetilde C_1(t-t')^{-2},
\]
and, for every $i\ge1$,
\[
 |\nabla_S^iH|^2\le\widetilde C_i(t-t')^{-(i+1)}.
\]
\end{theorem}
\begin{proof}
The curvature barrier gives the common $m=0$ estimate.  The graph-slope
lemma supplies uniform ellipticity and the ambient product coefficient bounds
are independent of $\lambda$.  Apply the parabolic regularity contract
\ref{lem:curve-flow-parabolic-regularity-interface} on the finite graph cover;
compactness of $P$ makes its constants simultaneous for all parameters.
The corrected dependence in Morgan--Tian's 2015 correction then gives the
stated dependence on $L_0$ and $\Theta_0$, and the estimates are uniform in
$0<\lambda<1$.
\end{proof}

\section{Controlled Curve-Shrinking Interface}

\begin{theorem}[Controlled curve shrinking for compact loop families]
\label{thm:controlled-curve-shrinking-existence}
\dcref{MT:ch19:19.1,MT:ch19:19.10,MT:ch19:19.14,MT:ch19:19.17,MT:ch19:19.24}
\source{morgan-tian}
\source{morgan-tian-correction}
\uses{thm:parametric-curve-shrinking-wellposedness,
  thm:ramp-curve-shrinking-global-existence,
  thm:uniform-controlled-ramp-approximation,
  lem:free-loop-approximation-preserves-pi1-orbit,
  thm:curve-flow-length-total-curvature-bounds,
  thm:curve-shrinking-derivative-estimates}
\notready
This is an assembled interface: the cited Morgan--Tian claims supply the
free-loop approximation, ramp preservation/global existence, and flow pieces,
while the correction supplies the required dependence on initial length and
total curvature.  The compact-family conclusion is the composition of the
explicit interfaces stated above.
Let $g(t)$ be a Ricci flow on a compact manifold for
$t\in[t_0,t_1]$, and let
$\Gamma\colon S^2\to\Lambda M$ be a continuous family of null-homotopic
loops.  Then for every $\zeta>0$ there is a homotopic family of smooth loops
$\widetilde\Gamma$ with length and least spanning area within $\zeta$ of
those of $\Gamma$, and constants $L_0,\Theta_0<\infty$ such that, for every
$c\in S^2$ and $0<\lambda<1$, the degree-one lift
$\widetilde\Gamma(c)^\lambda\subset M\times S^1_\lambda$ has a unique
geometric curve-shrinking flow, up to time-dependent reparametrization,
throughout $[t_0,t_1]$.  For each fixed $\lambda$, these flows form a
continuous family in $c$, remain ramps, obey uniform length and
total-curvature bounds depending
on $L_0,\Theta_0$ and the ambient flow, and satisfy the uniform local
curvature and derivative estimates above whenever the small-turning
hypothesis holds.
The family statement means continuity as a map
$S^2\times[t_0,t_1]\to\Lambda M$ after projection from the chosen lifted
branches in $M\times\mathbb R$.  This is a free-loop family statement; no
based representative is fixed, but the free-loop orbit is preserved by
\ref{lem:free-loop-approximation-preserves-pi1-orbit} whenever
$\pi_2(M)=0$.
\end{theorem}

\begin{proof}
Use \ref{thm:uniform-controlled-ramp-approximation} to choose
$\widetilde\Gamma$ and uniform initial bounds $L_0,\Theta_0$.  The lifts are
ramps, so \ref{thm:ramp-curve-shrinking-global-existence} gives a flow for
every parameter through $t_1$.  For fixed $\lambda$, the approximation family
is continuous in $C^{2,\alpha}$, and the circle component gives the uniform
initial immersion bound $|\partial_x\widetilde\Gamma(c)^\lambda|
\ge\lambda/(2\pi)$.  For each fixed $\lambda$,
\ref{thm:parametric-curve-shrinking-wellposedness} therefore supplies
uniqueness and continuous dependence on a common initial interval.  The
uniform ramp and curvature bounds allow continuation on
finitely many overlapping intervals covering $[t_0,t_1]$.  The estimates of
\ref{thm:curve-flow-length-total-curvature-bounds} are uniform because their
initial constants are $L_0,\Theta_0$; the local curvature and higher
derivative conclusions are exactly
\ref{thm:curve-shrinking-derivative-estimates} with those same uniform
constants.
\end{proof}
